\chapter{Behavioral Macroeconomics}
\label{chap:behavioral-macro}

Belief rules shared by the behavioral macro models: a linear rule $(m,\theta)$ maps today's and yesterday's rational
forecasts $e,e'$ to $me+\theta(e-e')$. Rational expectations is $(1,0)$, cognitive discounting is $(m,0)$ and
diagnostic expectations is $(1,\theta)$.

\section{Belief rules}

\begin{definition}[Linear belief rules]
\label{def:bm-belief-rule}
\lean{Econ.BeliefRule.Linear, Econ.BeliefRule.Linear.forecast, Econ.BeliefRule.Linear.rational, Econ.BeliefRule.Linear.cognitive, Econ.BeliefRule.Linear.diagnostic}
\leanok
$\hat E[y]=m\,e+\theta(e-e')$.
\end{definition}

\begin{lemma}[Rational expectations as a limit]
\label{lem:bm-belief-rule}
\uses{def:bm-belief-rule}
\lean{Econ.BeliefRule.Linear.forecast_rational, Econ.BeliefRule.Linear.forecast_cognitive, Econ.BeliefRule.Linear.forecast_diagnostic, Econ.BeliefRule.Linear.cognitive_one, Econ.BeliefRule.Linear.diagnostic_zero, Econ.BeliefRule.Linear.forecast_eq_iff, Econ.BeliefRule.Linear.continuous_forecast}
\leanok
Both behavioral families pass through the rational rule, which is the only rule forecasting $e$ for all inputs,
and forecasts are continuous in $(m,\theta)$.
\end{lemma}
\begin{proof}
\leanok
Evaluate at $(e,e')=(1,1)$ and $(1,0)$.
\end{proof}

\section{Cognitive discounting in the New Keynesian model}

\begin{definition}[Behavioral NK model]
\label{def:bm-cnk}
\uses{def:bm-belief-rule}
\lean{Econ.CognitiveNK.Model, Econ.CognitiveNK.Model.IsEquilibrium, Econ.CognitiveNK.Model.Determinate, Econ.CognitiveNK.Model.forward, Econ.CognitiveNK.Model.index, Econ.CognitiveNK.Model.toRational}
\leanok
$x_t=Mx_{t+1}-(i_t-\pi_{t+1})/\sigma$, $\pi_t=\beta M^f\pi_{t+1}+\kappa x_t$, $i_t=\phi_\pi\pi_t+\phi_x x_t$, with
$0<M,M^f\le1$. Determinacy means the only bounded equilibrium is zero.
\end{definition}

\begin{theorem}[Gabaix: determinacy with cognitive discounting]
\label{thm:bm-cnk-det}
\uses{def:bm-cnk}
\lean{Econ.CognitiveNK.Model.forward_det, Econ.CognitiveNK.Model.forward_trace, Econ.CognitiveNK.Model.isEquilibrium_iff, Econ.CognitiveNK.Model.determinate_iff_forward, Econ.CognitiveNK.Model.determinate_iff, Econ.CognitiveNK.Model.peg_determinate_iff, Econ.CognitiveNK.Model.peg_indeterminate_rational, Econ.CognitiveNK.Model.isEquilibrium_rational, Econ.CognitiveNK.Model.determinate_iff_rational, Econ.CognitiveNK.Model.determinate_robust}
\leanok
The model is determinate iff $\kappa(\phi_\pi-1)+(1-\beta M^f)\phi_x+\sigma(1-M)(1-\beta M^f)>0$. A peg is determinate
iff $\kappa<\sigma(1-M)(1-\beta M^f)$, never under rational expectations. At $M=M^f=1$ the model and its determinacy
condition are those of the rational NK model, and the rational Taylor condition implies determinacy for all
$M,M^f\in(0,1]$.
\end{theorem}
\begin{proof}
\leanok
Equilibria are orbits of the forward matrix, with $\det=(\sigma+\phi_x+\kappa\phi_\pi)/(\sigma M\beta M^f)>1$ and
$1+\det-\mathrm{tr}=\text{index}/(\sigma M\beta M^f)$; apply the two-dimensional determinacy criterion of the NK chapter.
\end{proof}

\begin{theorem}[Forward guidance: puzzle and resolution]
\label{thm:bm-fg}
\uses{def:bm-cnk}
\lean{Econ.CognitiveNK.Model.IsISPath, Econ.CognitiveNK.Model.cut, Econ.CognitiveNK.Model.fg_response, Econ.CognitiveNK.Model.fg_path_exists, Econ.CognitiveNK.Model.fg_vanishes, Econ.CognitiveNK.Model.fg_puzzle_rational, Econ.CognitiveNK.Model.fg_eventually_below_half}
\leanok
In partial equilibrium (the IS block with an exogenous real-rate path and output back at steady state after the cut), a unit real-rate cut at date $T$ followed by a return to steady state raises output today by $M^T/\sigma$.
For $M<1$ this vanishes as $T\to\infty$; at $M=1$ it is $1/\sigma$ at every horizon, so undiminished forward guidance
is knife-edge.
\end{theorem}
\begin{proof}
\leanok
Backward induction from $T$: $x_{T-k}=M^k/\sigma$.
\end{proof}

\section{Diagnostic expectations}

\begin{definition}[Forecasts in an AR(1) window]
\label{def:bm-de}
\uses{def:bm-belief-rule}
\lean{Econ.Diagnostic.Var, Econ.Diagnostic.cov, Econ.Diagnostic.xm2, Econ.Diagnostic.xm1, Econ.Diagnostic.x0, Econ.Diagnostic.x1, Econ.Diagnostic.εt, Econ.Diagnostic.fcT, Econ.Diagnostic.fcT1, Econ.Diagnostic.FE, Econ.Diagnostic.FR, Econ.Diagnostic.cgCoeff}
\leanok
Variables are coefficient vectors on the uncorrelated $(x_{t-2},\varepsilon_{t-1},\varepsilon_t,\varepsilon_{t+1})$;
forecasts of $x_{t+1}$ apply a linear belief rule to the rational forecasts; the CG coefficient is
$\mathrm{Cov}(FE,FR)/\mathrm{Var}(FR)$.
\end{definition}

\begin{theorem}[Over- and underreaction in survey forecasts]
\label{thm:bm-cg}
\uses{def:bm-de}
\lean{Econ.Diagnostic.FE_eq, Econ.Diagnostic.FR_eq, Econ.Diagnostic.cov_FE_FR, Econ.Diagnostic.cov_FR_FR, Econ.Diagnostic.cov_FE_FR_rational, Econ.Diagnostic.cov_FE_FR_diagnostic, Econ.Diagnostic.cov_FE_FR_cognitive, Econ.Diagnostic.cgCoeff_neg_iff, Econ.Diagnostic.cgCoeff_cognitive_pos, Econ.Diagnostic.fcT_eq, Econ.Diagnostic.fcT_diagnostic, Econ.Diagnostic.fe_orthogonal_iff}
\leanok
$\mathrm{Cov}(FE,FR)=s^2\rho^2[(1-m-\theta)(m+\theta)-\theta(1-m)\rho^2]$. Diagnostic expectations give
$-\theta(1+\theta)\rho^2s^2$, a negative CG coefficient iff $\theta>0$ (overreaction); cognitive discounting gives
$(1-m)m\rho^2s^2>0$ under horizon-independent shrinkage (underreaction); rational expectations give $0$. Forecast errors are uncorrelated with date-$t$
information iff the rule is rational.
\end{theorem}
\begin{proof}
\leanok
$FE=((1-m)\rho^3,(1-m)\rho^2,(1-m-\theta)\rho,1)$ and $FR=(0,-\theta\rho^2,(m+\theta)\rho,0)$.
\end{proof}

\begin{theorem}[Bordalo--Gennaioli--Shleifer: overreaction and reversal in prices]
\label{thm:bm-bgs-price}
\uses{def:bm-de}
\lean{Econ.Diagnostic.forecastIRF, Econ.Diagnostic.priceIRF, Econ.Diagnostic.hasSum_priceIRF, Econ.Diagnostic.priceIRF_eq, Econ.Diagnostic.price_overreaction, Econ.Diagnostic.price_reversal, Econ.Diagnostic.price_change_lt}
\leanok
With $|\delta\rho|<1$, the diagnostic price of a claim to $(x_{t+k})_{k\ge1}$ responds to a unit shock by $(1+\theta)c$
on impact and by $\rho^kc$ afterwards, $c=\delta\rho/(1-\delta\rho)$: overreaction followed by reversal, so the price
change after good news is predictably below its rational counterpart.
\end{theorem}
\begin{proof}
\leanok
Geometric series; the diagnostic term only loads on news, which is present on impact only.
\end{proof}

\section{Forward guidance in general equilibrium}

\begin{definition}[Forward-guidance experiment]
\label{def:bm-fgge}
\uses{def:bm-cnk}
\lean{Econ.CognitiveNK.iter, Econ.CognitiveNK.iter_bound, Econ.CognitiveNK.iter_tendsto_zero, Econ.CognitiveNK.Model.IsRatePath, Econ.CognitiveNK.Model.response}
\leanok
The nominal rate is at steady state except for a unit cut at date $T$, and the economy is at steady state after $T$.
Solving backward, $(x_0,\pi_0)=B^T(1/\sigma,\kappa/\sigma)$ with $B=\begin{pmatrix}M&1/\sigma\\ \kappa M&\beta M^f+\kappa/\sigma\end{pmatrix}$.
\end{definition}

\begin{theorem}[Gabaix: forward-guidance puzzle and its resolution]
\label{thm:bm-fgge}
\uses{def:bm-fgge, thm:bm-cnk-det}
\lean{Econ.CognitiveNK.Model.fgGE_response, Econ.CognitiveNK.Model.fgGE_path_exists, Econ.CognitiveNK.Model.fgGE_puzzle_bound, Econ.CognitiveNK.Model.fgGE_puzzle, Econ.CognitiveNK.Model.fgGE_vanishes, Econ.CognitiveNK.Model.iter_nonneg, Econ.CognitiveNK.Model.fgGE_not_vanishes, Econ.CognitiveNK.Model.fgGE_vanishes_iff, Econ.CognitiveNK.Model.fgGE_vanishes_iff_determinate}
\leanok
Under rational expectations the output response to a cut $T$ periods ahead is at least $1/\sigma+T\kappa/\sigma^2$, so
it grows without bound in $T$. If cognitive discounting makes a peg determinate, $\kappa<\sigma(1-M)(1-\beta M^f)$,
the output and inflation responses tend to $0$ as $T\to\infty$; conversely the output response vanishes only if the
peg is determinate. (Under rational expectations the return-to-steady-state path is a selection among bounded solutions.)
\end{theorem}
\begin{proof}
\leanok
At $M=M^f=1$, $\pi$ stays above $\kappa/\sigma$ and each backward step adds at least $\kappa/\sigma^2$ to $x$.
Under the peg condition, $w=(1,w_2)$ with $w_2$ between $\kappa M/(1-\beta M^f-\kappa/\sigma)$ and $\sigma(1-M)$ satisfies
$Bw\le\rho w$ for some $\rho<1$, so the nonnegative iterates are bounded by $c\rho^T w$. Otherwise
$\ell=(1,(1-M)/(\kappa M))$ satisfies $\ell B\ge\ell$, so $\ell\cdot(x_T,\pi_T)$ stays above $1/\sigma$.
\end{proof}

\section{Diagnostic credit cycles}

\begin{theorem}[Bordalo--Gennaioli--Shleifer: credit cycles]
\label{thm:bm-credit}
\uses{def:bm-de, thm:bm-bgs-price}
\lean{Econ.CreditCycles.spreadIRF, Econ.CreditCycles.investIRF, Econ.CreditCycles.forecastIRF_diagnostic, Econ.CreditCycles.spreadIRF_eq, Econ.CreditCycles.investIRF_eq, Econ.CreditCycles.spread_too_low, Econ.CreditCycles.invest_boom, Econ.CreditCycles.spread_reversal, Econ.CreditCycles.invest_bust, Econ.CreditCycles.invest_bust_lt, Econ.CreditCycles.spread_reversal_pos, Econ.CreditCycles.spread, Econ.CreditCycles.fcNext, Econ.CreditCycles.spreadNext, Econ.CreditCycles.condExp, Econ.CreditCycles.fcNext_eq, Econ.CreditCycles.condExp_spreadNext, Econ.CreditCycles.spread_sentiment, Econ.CreditCycles.expected_spread_change_eq, Econ.CreditCycles.cov_expected_change}
\leanok
With spreads $S_t=-\chi\hat E_ta_{t+1}$ and investment $k_{t+1}=\hat E_ta_{t+1}/(1-\alpha)$, after good news spreads fall and
investment rises by more than under rational expectations; next period spreads rise by $\chi\theta\rho$ more and
investment falls by $\theta\rho/(1-\alpha)$ more. The expected future spread $E_tS_{t+1}=-\chi\rho^2a_t$ does not depend on $\theta$, so the
spread's sentiment component $-\chi\theta\rho\varepsilon_t$ fully reverts: the expected spread change exceeds its rational
value by $\chi\theta\rho\varepsilon_t$, which covaries positively with news and negatively with sentiment.
\end{theorem}
\begin{proof}
\leanok
The diagnostic forecast response is $(1+\theta[k=0])\rho^{k+1}$, and $\hat E_ta_{t+1}=\rho a_t+\theta\rho\varepsilon_t$.
\end{proof}
